%% kapa: exam=2 round=27 year=2016 period=3 subject=compensation_law subject_ko=감정평가 및 보상법규 seq=2 number=2 points=30
%% source: Q-Net 감정평가사 기출문제 (cst00309) artlSeq=5185401 file=extracted/1008/5185401/제27회 감정평가사 2차시험 기출문제/제27회 감정평가사 2차 시험문제지/(3교시 시험문제지)감정평가 및 보상법규.pdf
%% conversion: kordoc 4.24.51 -> tools/gen2.py (auto); review=proofread-gpt-sol; flags=none
\begin{longquestion}{【문제 2】}
甲은 2015. 3. 16. 乙로부터 A광역시 B구 소재 도로로 사용되고 있는 토지 200m\textsuperscript{2}(이하 ‘이 사건 토지’라 함)를 매수한 후 자신의 명의로 소유권 이전등기를 하였다. 한편, 甲은 A광역시지방토지수용위원회에 “사업시행자인 B구청장이 도로개설공사를 시행하면서 사업인정고시가 된 2010. 4. 6. 이후 3년 이상 이 사건 토지를 사용하였다”고 주장하면서 「공익사업을 위한 토지 등의 취득 및 보상에 관한 법률」(이하 ‘토지보상법’이라 함) 제72조제1호를 근거로 이 사건 토지의 수용을 청구하였다. 이에 대해 A광역시지방토지수용위원회는 “사업인정고시가 된 날부터 1년 이내에 B구청장이 재결신청을 하지 아니하여 사업인정은 그 효력을 상실하였으므로 甲은 토지수용법 제72조제1호를 근거로 이 사건 토지의 수용을 청구할 수 없다”며 甲의 수용청구를 각하하는 재결을 하였다. 다음 물음에 답하시오. (30점) \par
1) A광역시지방토지수용위원회의 각하재결에 대하여 행정소송을 제기하기 전에 강구할 수 있는 甲의 권리구제수단에 관하여 설명하시오. (10점) \par
2) 甲이 A광역시지방토지수용위원회의 각하재결에 대하여 행정소송을 제기할 경우 그 소송의 형태와 피고적격에 관하여 설명하시오. (20점) \par
% [원문 그대로] 본문의 ‘토지수용법 제72조제1호’는 원문 표기를 보존함.
\begin{framedtext}
참조 조문 \par
<공익사업을 위한 토지 등의 취득 및 보상에 관한 법률> \par
제23조(사업인정의 실효) ① 사업시행자가 제22조제1항에 따른 사업인정의 고시(이하 "사업인정고시"라 한다)가 된 날부터 1년 이내에 제28조제1항에 따른 재결신청을 하지 아니한 경우에는 사업인정고시가 된 날부터 1년이 되는 날의 다음 날에 사업인정은 그 효력을 상실한다. \par
② <이하 생략> \par
제72조(사용하는 토지의 매수청구 등) 사업인정고시가 된 후 다음 각 호의 어느 하나에 해당할 때에는 해당 토지소유자는 사업시행자에게 해당 토지의 매수를 청구하거나 관할 토지수용위원회에 그 토지의 수용을 청구할 수 있다. 이 경우 관계인은 사업시행자나 관할 토지수용위원회에 그 권리의 존속(存續)을 청구할 수 있다. \par
1. 토지를 사용하는 기간이 3년 이상인 경우 \par
2. <이하 생략> \par
\end{framedtext}
\end{longquestion}
